% ======================================================================
%  sec14.tex —— 第 14 节 拟凝聚代数
%  底本：原书 pp.72–76（印刷页 = PDF 页）
%  原书标题：14. Quasi-coherent algebras
%  体例：\section + \subsection 与原文编号一致；定理类环境紧跟 \label{类型:编号}。
% ======================================================================

\section{拟凝聚代数}\label{sec:14}

% ---------------------------------------------------------------- p.72
\subsection{定义}\label{sec:14.1}

\begin{definition}\label{def:14.1.1}
设 $(X,\cO_X)$ 是环化空间。$\cO_X$-\emph{代数}（层）是指一个对象 $\cA\in\Shv(X,\CRing)$，连同态射 $\cO_X\to\cA$。记
\[
\cO_X\text{-}\mathsf{alg}:=\Shv(X,\CRing)_{\cO_X/}
\]
为 $\cO_X$-代数构成的范畴。
\end{definition}

\begin{construction}\label{con:14.1.2}
设 $X$ 是概形，$\cI\subseteq\cO_X$ 是理想层。考虑 $\cO_X$-模
\[
\cO_X/\cI:=\coker(\cI\to\cO_X)\in\omod.
\]
展开定义可知，它同构于 $\cO_X$-模预层 $(\cO_X/\cI)_{\PShv}$ 的层化，后者把 $U$ 送到 $\cO_X(U)/\cI(U)$。注意 $\cO_X(U)/\cI(U)$ 带有显然的 $\cO_X(U)$-代数结构。由此推出 $\cO_X/\cI$ 继承 $\cO_X$-代数（层）结构。
\end{construction}

\begin{construction}\label{con:14.1.3}
有遗忘函子
\[
\cO_X\text{-}\mathsf{alg}\to\omod,
\]
它把 $\cO_X$-代数 $\cA$ 送到底层的阿贝尔层 $\cA$，并配上由
\[
\cO_X(U)\times\cA(U)\to\cA(U)\times\cA(U)\to\cA(U)
\]
给出的 $\cO_X$-模结构。
\end{construction}

\begin{definition}\label{def:14.1.4}
设 $X$ 是概形。称 $\cO_X$-代数 $\cA$ \emph{拟凝聚}，若把它视为 $\cO_X$-模时拟凝聚。记
\[
\cO_X\text{-}\mathsf{alg}_{\mathrm{qcoh}}\subseteq\cO_X\text{-}\mathsf{alg}
\]
为拟凝聚 $\cO_X$-代数构成的全子范畴。
\end{definition}

\begin{definition}\label{def:14.1.5}
设 $X$ 是概形。$\cO_X$ 的\emph{拟凝聚理想}是指这样的理想层 $\cI$：把它视为 $\cO_X$-模时拟凝聚。
\end{definition}

\begin{example}\label{ex:14.1.6}
设 $X$ 是概形，$\cI\subseteq\cO_X$ 是拟凝聚理想。则 $\cO_X/\cI$ 是拟凝聚 $\cO_X$-代数。
\end{example}

\begin{example}\label{ex:14.1.7}
设 $A\in\CRing$，$X:=\Spec(A)$。对任意 $B\in A\text{-}\mathsf{alg}$，可以把 $B$ 视为 $A$-模，并考虑拟凝聚 $\cO_X$-模 $\widetilde{B}$。回忆 $\widetilde{B}$ 是 $\cO_X$-模预层
\[
U\mapsto\cO_X(U)\otimes_AB
\]
的层化。注意 $\cO_X(U)\otimes_AB$ 带有显然的 $\cO_X(U)$-代数结构。由此推出 $\widetilde{B}$ 继承 $\cO_X$-代数（层）结构。
\end{example}

\begin{exercise}\label{exc:14.1.8}
设 $A\in\CRing$，$X:=\Spec(A)$。证明下述函子互为逆函子：
\[
\begin{array}{ccc}
A\text{-}\mathsf{alg} & \rightleftarrows & \cO_X\text{-}\mathsf{alg}_{\mathrm{qcoh}}\\[0.5ex]
B & \mapsto & \widetilde{B}\\
\cB(X) & \mapsfrom & \cB.
\end{array}
\]
\end{exercise}

\begin{yizhu}
习题 \ref{exc:14.1.8} 的等价在 Stacks 同一节里还有两条常用换算：$\widetilde{M\otimes_RN}\simeq\widetilde{M}\otimes_{\cO_X}\widetilde{N}$，以及 $\widetilde{\mathrm{Sym}^nM}\simeq\mathrm{Sym}^n\widetilde{M}$。后者正是 \ref{sec:14.3} 的代数基础：先取对称代数、再取相对谱，两步都在仿射开集上逐点计算。
[来源：Stacks Project, Tag 01I6（Section 26.7, Quasi-coherent sheaves on affines），\url{https://stacks.math.columbia.edu/tag/01I6}（2026-11-17 访问）。]
\end{yizhu}

% ---------------------------------------------------------------- p.73
\subsection{相对 Spec 构造}\label{sec:14.2}

\begin{construction}\label{con:14.2.1}
设 $S$ 是概形，$p\colon X\to S$ 是 $S$-概形。对象 $p_*\cO_X\in\Shv(S,\CRing)$ 带有 $\cO_S$-代数结构，该结构由典范态射 $\cO_X\to p_*\cO_X$ 给出。
\end{construction}

\begin{example}\label{ex:14.2.2}
设 $p\colon X\to S$ 是概形之间的闭浸入，$\cI_X\subseteq\cO_S$ 是其定义理想。回忆有阿贝尔层之间的同构
\[
\cO_S/\cI_X\simeq p_*\cO_X.
\]
容易验证，它也是 $\cO_S$-代数之间的同构。注意由引理 \ref{lem:13.4.4}，两者都拟凝聚。特别地，
\[
\cI_X\simeq\ker(\cO_S\to p_*\cO_X)
\]
是 $\cO_S$ 的拟凝聚理想。
\end{example}

\begin{construction}\label{con:14.2.3}
设 $S$ 是概形，$p\colon X\to S$、$q\colon Y\to S$ 是 $S$-概形。设 $f\colon X\to Y$ 是 $S$-态射，即满足 $p=q\circ f$ 的态射。有典范的 $\cO_S$-同态
\[
q_*\cO_Y\to q_*(f_*\cO_X)\simeq p_*\cO_X.
\]
容易验证，这定义了反变函子
\[
\begin{array}{rcl}
(\Sch_S)^{\op} & \to & \cO_S\text{-}\mathsf{alg}\\
(X\xrightarrow{p}S) & \mapsto & p_*\cO_X.
\end{array}
\]
\end{construction}

下述定理说明，该函子部分定义的左伴随在拟凝聚 $\cO_S$-代数上有定义。

\begin{theorem}\label{thm:14.2.4}
设 $S$ 是概形。对任意 $\cA\in\cO_S\text{-}\mathsf{alg}_{\mathrm{qcoh}}$，存在本质唯一的 $S$-概形 $q\colon Y\to S$，连同同构 $\cA\xrightarrow{\simeq}q_*\cO_Y$，使得对任意满足 $p_*\cO_X$ 拟凝聚的测试 $S$-概形 $p\colon X\to S$，复合
\begin{equation}\tag{14.1}\label{eq:14.1}
\Hom_{\Sch_S}(X,Y)\to\Hom_{\cO_S\text{-}\mathsf{alg}}(q_*\cO_Y,p_*\cO_X)\xrightarrow{\simeq}\Hom_{\cO_S\text{-}\mathsf{alg}}(\cA,p_*\cO_X)
\end{equation}
是双射。把这个 $S$-概形记作
\[
\Spec_S(\cA)\in\Sch_S.
\]
\end{theorem}

\begin{yizhu}
Stacks 把相对谱的泛性质写成「$\Hom_{\Sch_S}(T,\underline{\Spec}_S(\cA))$ 与满足 $\varphi\colon f^*\cA\to\cO_T$ 的配对 $(f\colon T\to S,\varphi)$ 之间的双射」。由伴随 $f^*\dashv f_*$，$\varphi$ 与 $\cA\to f_*\cO_T$ 互相转换；所以这与 \eqref{eq:14.1} 里的 $\Hom_{\cO_S\text{-}\mathsf{alg}}(\cA,p_*\cO_X)$ 是同一件事，只是把同构放宽成了任意态射。
[来源：Stacks Project, Tag 01LQ（Section 27.4, Relative spectrum as a functor），\url{https://stacks.math.columbia.edu/tag/01LQ}（2026-11-17 访问）。]
\end{yizhu}

\begin{sketch}
（本质）唯一性由 Yoneda 引理推出，只需证明存在性。

先设 $S=\Spec(R)$ 仿射。令 $A:=\cA(S)$ 是对应于 $\cA$ 的 $R$-代数，写 $Y:=\Spec(A)$，并通过对应于 $R\to A$ 的态射 $q\colon\Spec(A)\to\Spec(R)$ 把 $Y$ 视为 $S$-概形。由习题 \ref{exc:14.1.8}，有 $\cA\simeq q_*\cO_Y$，因为 $\cA(S)\simeq A\simeq(q_*\cO_Y)(S)$。

设 $p\colon X\to S$ 是使 $p_*\cO_X$ 拟凝聚的测试 $S$-概形，需验证 \eqref{eq:14.1} 是双射。设 $\phi\colon R\simeq\cO_S(S)\to\cO_X(X)$ 是由 $p$ 诱导的同态。展开定义，有
\[
\begin{aligned}
\Hom_{\Sch_S}(X,Y)&\simeq\Hom_{\Sch}(X,Y)\times_{\Hom_{\Sch}(X,S)}\{p\}\\
&\simeq\Hom_{\CRing}(A,\cO_X(X))\times_{\Hom_{\CRing}(R,\cO_X(X))}\{\phi\}\\
&\simeq\Hom_{R\text{-}\mathsf{alg}}(A,\cO_X(X))\\
&\simeq\Hom_{\cO_S\text{-}\mathsf{alg}_{\mathrm{qcoh}}}(\cA,p_*\cO_X),
\end{aligned}
\]
其中第二个双射由定理 7.1.2 给出，第四个由习题 \ref{exc:14.1.8} 给出。容易验证上述复合等于 \eqref{eq:14.1}，故后者也是双射。

% ---------------------------------------------------------------- p.74
现设 $S$ 是一般概形。取 $S$ 的仿射开子概形覆盖 $S=\bigcup_{i\in I}S_i$，并记 $S_{ij}:=S_i\cap S_j$。注意 $\cA_i:=\cA|_{S_i}$ 是拟凝聚 $\cO_{S_i}$-代数。由上述讨论，存在本质唯一的 $S_i$-概形 $q_i\colon Y_i\to S_i$，连同同构 $\cA_i\simeq q_{i,*}\cO_{Y_i}$，使得对任意满足 $p_*\cO_X$ 拟凝聚的 $S_i$-概形 $p\colon X\to S_i$，有
\[
\Hom_{\Sch_{S_i}}(X,Y_i)\simeq\Hom_{\cO_{S_i}\text{-}\mathsf{alg}}(\cA_i,p_*\cO_X).
\]
对 $(i,j)\in I^2$，考虑 $S_{ij}$-概形
\[
q_{ij}\colon Y_{ij}:=Y_i\times_{S_i}S_{ij}\to S_{ij}.
\]
容易验证
\[
q_{ij,*}\cO_{Y_{ij}}\simeq(q_{i,*}\cO_{Y_i})|_{S_{ij}}\simeq\cA|_{S_{ij}}.
\]
特别地，$q_{ij,*}\cO_{Y_{ij}}$ 拟凝聚。

对任意测试 $S_{ij}$-概形 $p\colon X\to S_{ij}$，有
\[
\Hom_{\Sch_{S_{ij}}}(X,Y_{ij})\simeq\Hom_{\Sch_{S_i}}(X,Y_i)\simeq\Hom_{\cO_{S_i}\text{-}\mathsf{alg}}(\cA_i,p'_*\cO_X),
\]
其中 $p'$ 是复合 $X\xrightarrow{p}S_{ij}\xrightarrow{u_{ij}}S_i$。又有
\[
\Hom_{\cO_{S_i}\text{-}\mathsf{alg}}(\cA_i,p'_*\cO_X)\simeq\Hom_{\cO_{S_i}\text{-}\mathsf{alg}}(\cA_i,s_{ij,*}\circ p_*\cO_X)\simeq\Hom_{\cO_{S_{ij}}\text{-}\mathsf{alg}}(\cA|_{S_{ij}},p_*\cO_X).
\]
这说明
\[
\Hom_{\Sch_{S_{ij}}}(X,Y_{ij})\simeq\Hom_{\cO_{S_{ij}}\text{-}\mathsf{alg}}(\cA|_{S_{ij}},p_*\cO_X).
\]
注意右端在交换 $i$ 与 $j$ 之后不变。于是得到典范双射
\[
\Hom_{\Sch_{S_{ij}}}(X,Y_{ij})\simeq\Hom_{\Sch_{S_{ji}}}(X,Y_{ji}).
\]
容易验证该双射对 $X$ 函子化。故由 Yoneda 引理，得到定义在 $S_{ij}=S_{ji}$ 上的同构
\[
\phi_{ij}\colon Y_{ij}\to Y_{ji}.
\]
容易验证
\[
(I,(Y_i)_{i\in I},(Y_{ij})_{(i,j)\in I^2},(\phi_{ij})_{(i,j)\in I^2})
\]
是概形的粘接数据（见定义 6.1.1）。设
\[
(Y,(Y'_i)_{i\in I},(\varphi_i)_{i\in I})
\]
是它的粘接输出（见命题–定义 6.1.2）。由命题 6.1.3，有典范态射
\[
q\colon Y\to S
\]
由 $Y_i\xrightarrow{q_i}S_i\to S$ 粘成。按构造，有同构
\[
\alpha_i\colon\cA|_{S_i}\simeq q_{i,*}\cO_{Y_i}\simeq q_*\cO_Y|_{S_i}.
\]
又 $\alpha_i$ 在 $S_{ij}$ 上的限制等于 $\alpha_j$ 在 $S_{ji}=S_{ij}$ 上的限制。故由习题 1.2.7，存在唯一同构
\[
\alpha\colon\cA\simeq q_*\cO_Y,
\]
使其在每个 $S_i$ 上的限制等于 $\alpha_i$。

对如上构造的 $q\colon Y\to S$ 与同构 $\cA\simeq q_*\cO_Y$，由命题–构造 6.1.4 可以验证：只要 $p_*\cO_X$ 拟凝聚，\eqref{eq:14.1} 就是双射。
\end{sketch}

\begin{corollary}\label{cor:14.2.5}
设 $S$ 是概形，$\Sch'_S\subseteq\Sch_S$ 是由满足 $p_*\cO_X$ 拟凝聚的 $S$-概形 $p\colon X\to S$ 构成的全子范畴。则函子
\[
(\Sch'_S)^{\op}\to\cO_S\text{-}\mathsf{alg}_{\mathrm{qcoh}},\qquad (p\colon X\to S)\mapsto p_*\cO_X
\]
有全忠实的左伴随
\[
\Spec_S\colon\cO_S\text{-}\mathsf{alg}_{\mathrm{qcoh}}\to(\Sch'_S)^{\op}.
\]
\end{corollary}

\begin{yizhu}
相对谱的存在性在 Stacks 中被单独陈述为「存在 $\pi\colon\underline{\Spec}_S(\cA)\to S$」，泛性质则另立一节。这样拆开的好处是：粘接构造只在存在性里出现一次，而一切用到相对谱的论证都只依赖泛性质。本节 Sketch 的处理顺序恰好相反，先对一般 $S$ 作粘接，再用泛性质收尾。
[来源：Stacks Project, Tag 01LP（Lemma 27.3.4），\url{https://stacks.math.columbia.edu/tag/01LP}（2026-11-17 访问）。]
\end{yizhu}

\begin{exercise}\label{exc:14.2.6}
设 $f\colon S'\to S$ 是概形之间的态射。对 $\cA\in\cO_S\text{-}\mathsf{alg}_{\mathrm{qcoh}}$，考虑
\[
\cA':=f^*\cA\in\cO_{S'}\text{-}\mathsf{alg}_{\mathrm{qcoh}}.
\]
证明
\[
\Spec_S(\cA)\times_SS'\simeq\Spec_{S'}(\cA').
\]
\end{exercise}

\begin{exercise}\label{exc:14.2.7}
设 $S$ 是概形，$\cA\in\cO_S\text{-}\mathsf{alg}_{\mathrm{qcoh}}$。证明对任意 $S$-概形 $p\colon X\to S$，典范映射
\[
\Hom_{\Sch_S}(X,\Spec_S(\cA))\to\Hom_{\cO_S\text{-}\mathsf{alg}}(\cA,p_*\cO_X)
\]
是双射。提示：沿 $X\to S$ 作基变换。
\end{exercise}

% ---------------------------------------------------------------- p.75
\subsection{向量丛}\label{sec:14.3}

\begin{construction}\label{con:14.3.1}
设 $(S,\cO_S)$ 是环化空间，$\cE$ 是 $\cO_S$-模。$\cE$ 的\emph{对称代数} $\mathrm{Sym}_{\cO_S}(\cE)$ 定义为预层
\[
U\mapsto\mathrm{Sym}_{\cO_S(U)}(\cE(U))
\]
的层化。
\end{construction}

下述结果的证明留给读者。

\begin{lemma}\label{lem:14.3.2}
设 $S$ 是概形，$\cE$ 是拟凝聚 $\cO_S$-模。则 $\mathrm{Sym}_{\cO_S}(\cE)$ 是拟凝聚 $\cO_S$-代数。此外，对任意仿射开子概形 $U\subseteq S$，有
\[
(\mathrm{Sym}_{\cO_S}(\cE))(U)\simeq\mathrm{Sym}_{\cO_S(U)}(\cE(U)).
\]
\end{lemma}

\begin{yizhu}
Stacks 的定义与此一致：$V(\cE)=\underline{\Spec}_S(\mathrm{Sym}(\cE))$，并指出 $\mathrm{Sym}(\cE)$ 拟凝聚，这正是引理 \ref{lem:14.3.2} 的第一句。该处还把 $V(\cE)$ 刻画成「$\cE$ 的对偶在测试概形上的整体截面」所给出的函子，与引理 \ref{lem:14.3.5} 是同一条双射的两种写法。
[来源：Stacks Project, Tag 01M1（Definition 27.6.1），\url{https://stacks.math.columbia.edu/tag/01M1}（2026-11-17 访问）。]
\end{yizhu}

\begin{construction}\label{con:14.3.3}
设 $S$ 是概形，$\cE$ 是拟凝聚 $\cO_S$-模。称
\[
V_S(\cE):=\Spec_S(\mathrm{Sym}_{\cO_S}(\cE))
\]
为与 $\cE$ 相伴的\emph{向量丛}。
\end{construction}

% ---------------------------------------------------------------- p.76
\begin{example}\label{ex:14.3.4}
有
\[
V_S(\cO_S^{\oplus n})\simeq\A^n_S:=\A^n_\Z\times S.
\]
特别地，若 $\cE$ 秩为 $n$ 局部自由，则 $S$-概形 $V_S(\cE)$ 与 $\A^n_S$ 在基变换到 $S$ 的某个开覆盖之后同构。
\end{example}

\begin{yizhu}
「秩 $n$ 局部自由」的准确含义是：每点有邻域 $U$，使 $\cE|_U\simeq\bigoplus_{i=1}^n\cO_S|_U$；一般情形只要求 $\cE|_U$ 是某个集合 $I$ 上的自由模，$I$ 有限时才称有限局部自由。须先作基变换才同构，说明 $V_S(\cE)$ 只是局部平凡，未必整体平凡，而「向量丛」一词正是对这种局部平凡性的概括。
[来源：Stacks Project, Tag 01C6（Definition 17.14.1），\url{https://stacks.math.columbia.edu/tag/01C6}（2026-11-17 访问）。]
\end{yizhu}

由习题 \ref{exc:14.2.7}，有下述结果。

\begin{lemma}\label{lem:14.3.5}
设 $S$ 是概形，$\cE$ 是拟凝聚 $\cO_S$-模。对任意 $S$-概形 $p\colon X\to S$，典范映射
\[
\Hom_{\Sch_S}(X,V_S(\cE))\to\Hom_{\cO_S\text{-}\mathsf{mod}}(\cE,p_*\cO_X)
\]
是双射。
\end{lemma}

\begin{yizhu}
取 $X=S$，引理 \ref{lem:14.3.5} 给出 $\Hom_{\Sch_S}(S,V_S(\cE))\simeq\Hom_{\cO_S}(\cE,\cO_S)$：$V_S(\cE)$ 的 $S$-点恰是 $\cE$ 的对偶模的整体截面。取 $\cE=\cO_S^{\oplus n}$ 便回到「$\A^n_S$ 的 $S$-点是 $n$ 个整体函数」这一熟悉事实，可见该双射把几何与代数两侧的「截面」对齐了。
[译者按（上下文推断，无外部依据）：这里对偶的出现源于 $\Spec_S(\mathrm{Sym}\,\cE)$ 的函数环是 $\mathrm{Sym}\,\cE$，故其点对应 $\cE$ 上的线性函数。]
\end{yizhu}
